The subset analyses above tested whether a \emph{group} of channels was sufficient as a self-contained detector; this analysis instead isolates what the remaining montage contributes to a single, fixed electrode. Each electrode was scored twice on the same sessions: once as the sole input to the complete pipeline (\emph{single electrode}), and once as that electrode's output within the full 32-channel Proposed-Med run, where Stage A screens epochs using every electrode (\emph{full montage}) (Table~\ref{tab:exp1_solo_vs_montage}). The single-electrode score requires no channel selection, so it is unaffected by the best-channel-per-session oracle used elsewhere; the paired difference, $\Delta F_1 = F_1^{\text{full montage}} - F_1^{\text{single electrode}}$, therefore isolates the gain contributed by the other 31 electrodes through Stage A's cross-channel epoch screening, with the target electrode held fixed.

At Fp1 and Fp2, this gain was not significant on either corpus, indicating that these two electrodes already capture nearly all of the detectable blink signal on their own: in Raja, $\Delta F_1=+0.0259$ for Fp1 ($p=1.000$) and $+0.0345$ for Fp2 ($p=1.000$); in Cao2018, $\Delta F_1=+0.0203$ for Fp1 ($p=1.000$) and $+0.0286$ for Fp2 ($p=0.991$). At every weaker frontal-adjacent electrode, by contrast, the gain was significant, showing that these electrodes rely on the rest of the montage: in Raja, $\Delta F_1=+0.0516$ for AF4 ($p=0.001$), $+0.0617$ for AF3 ($p=0.002$), $+0.0579$ for F3 ($p=0.008$), and $+0.0526$ for F4 ($p<0.001$); in Cao2018, $\Delta F_1=+0.1021$ for F3 ($p<0.001$) and $+0.0896$ for F4 ($p<0.001$). For these weaker electrodes, the gain was concentrated in recall rather than precision, consistent with the full montage recovering blink epochs that Stage A misses when screening on a single noisy channel.
